\chapter{Les environnements logiciels}
\label{chap:D-environnements}

Cette annexe donne les versions des logiciels qui ont produit les résultats, la plateforme et la
chaîne qui construit les documents. Les fichiers de \file{envs/} font foi. Le \cref{tab:env} en
extrait les paquets clés.

\section{La plateforme}
\label{sec:d-plateforme}

Tous les essais et toutes les analyses publiées viennent d'un Mac à puce Apple M2 Max, sous macOS 27.2,
pour l'architecture arm64, avec Python 3.11.18. Les simulations ont tourné sur le processeur central,
un fil d'exécution par essai. Les essais GRPO ont tourné sur le processeur graphique de la puce
(interface MPS de PyTorch) ou sur le processeur central, selon l'essai
(\cref{chap:07-modeles-de-langue}). Les environnements d'origine viennent de conda
(\texttt{conda create -p sim/env python=3.11}), puis pip a installé chaque paquet. Un environnement
\texttt{venv} ordinaire reçoit les mêmes paquets.

\section{Les quatre environnements}
\label{sec:d-envs}

Chaque fichier de \file{envs/} est le \texttt{pip freeze} d'un environnement, avec la
version de Python et la plateforme en tête. Les fichiers \file{*.lock} donnent les mêmes listes avec
l'empreinte SHA-256 de chaque paquet, pour une installation par \texttt{pip install
-{}-require-hashes}. Quatre fichiers \file{*.txt}, un par environnement, reprennent les listes que
citent les protocoles, \file{sim/requirements.lock} et ses voisins, réunies en un seul dossier. Le
cinquième, \file{test.txt}, ajoute les paquets des tests à l'environnement d'analyse.

\begin{table}[htbp]
\centering
\small
\caption{Les environnements figés du projet (\file{envs/README.md}).}
\label{tab:d-envs}
\begin{tabularx}{\linewidth}{@{}llX@{}}
\toprule
Fichier & Environnement & Ce qu'il fait tourner \\
\midrule
\file{train.txt} & \file{sim/env} & les entraînements (\file{sim/run.py}) et les réévaluations
  (\file{sim/dose.py}) \\
\file{analysis.txt} & \file{sim/env-an} & chaque script d'analyse, \file{docs/check\_numbers.py},
  l'intégration continue \\
\file{zoo.txt} & \file{sim/env-zoo} & le pipeline officiel de RL Baselines3 Zoo de l'expérience E6 \\
\file{llm.txt} & \file{sim/env-llm} & GRPO (\file{sim/llm/grpo.py}) et les figures de \file{docs/} \\
\file{test.txt} & sur \file{analysis.txt} & les tests unitaires (\file{tests/}) \\
\bottomrule
\end{tabularx}
\end{table}

\gentable{env}

\section{Les pièges de l'installation}
\label{sec:d-installer}

\begin{itemize}
  \item \textbf{Recréer un environnement.} Les chemins comptent, puisque les scripts et les
        lanceurs appellent \file{sim/env/bin/python} et ses voisins :
\begin{lstlisting}[language=bash]
# Créer l'environnement d'analyse
python3.11 -m venv sim/env-an
# Installer les paquets de l'analyse et des tests
sim/env-an/bin/pip install -r envs/analysis.txt -r envs/test.txt
# Tout recalculer et tout vérifier
make verify
\end{lstlisting}
        Les trois autres suivent le même modèle. \file{sim/env-zoo} est un \texttt{venv} construit
        sur le Python de \file{sim/env}.
  \item \textbf{rliable 1.2.0} n'est distribué que sous forme d'archive source. pip le construit en
        quelques secondes : c'est du Python pur.
  \item \textbf{Sur Linux}, \file{analysis.txt} s'installe tel quel. L'intégration continue le fait
        sur Ubuntu. Les derniers chiffres d'un bootstrap bougent avec le processeur : la comparaison
        octet par octet est exigée sur macOS arm64, la plateforme des résultats. Linux est
        comparé à $10^{-12}$ près en relatif (\cref{sec:ci}).
  \item \textbf{Entraîner sur Linux} fonctionne avec \file{train.txt}, mais les essais du dépôt
        viennent de macOS arm64. Un nouvel essai ailleurs donne d'autres courbes pour la même
        graine aléatoire : MuJoCo, JAX et PyTorch arrondissent différemment d'un processeur à l'autre
        (\cref{sec:reentrainer}).
  \item \textbf{GRPO} télécharge \texttt{Qwen/Qwen2.5-0.5B-Instruct} et
        \texttt{unsloth/Llama-3.2-1B-Instruct} depuis Hugging Face. Depuis l'écart 15 (c) du volet
        GRPO, \file{grpo.py} fixe la révision exacte de chaque modèle, celle de nos essais.
                \file{docs/figures/build.py} télécharge le découpeur en jetons de Qwen. Le fichier
        \file{generation\_config.json} de Qwen fixe une pénalité de répétition de 1,1, que les
        essais publiés ont subie (écart 18 du volet GRPO). L'option \texttt{-{}-c7} de \file{grpo.py}
        la ramène à 1,0 pour la relance de l'écart 20. Sans l'option, le code refait les anciens essais à l'identique.
\end{itemize}

\section{La chaîne des documents}
\label{sec:d-chaine}

Hors des environnements Python, la construction des PDF et des figures emploie les outils du
\cref{tab:d-chaine}. Ces outils dessinent et composent sans calculer aucun nombre.
\texttt{make check} et \texttt{make verify} n'ont besoin que de \file{analysis.txt}.

\begin{table}[htbp]
\centering
\small
\caption{Les outils de construction des documents.}
\label{tab:d-chaine}
\begin{tabularx}{\linewidth}{@{}llX@{}}
\toprule
Outil & Version & Employé par \\
\midrule
pandoc & 3.6.1 & \file{docs/build.sh} \\
md2star & 3.13.5 (PyPI) & \file{docs/build.sh} (DOCX et un premier PDF) \\
TeX Live 2026 : LuaHBTeX, latexmk & 1.24.0, 4.88 & \file{docs/build.sh}, \file{report/build.sh} \\
sprezzature-figures & 2.14.1 (2.14.0 sur PyPI) & \file{docs/figures/build.py},
  \file{article/figures/build.py} \\
rsvg-convert (librsvg) & 2.61.3 & conversion des SVG en PDF et en PNG \\
ffmpeg & 8.1.2 & \file{report/tools/new\_figures.py} (images tirées des vidéos de E6) \\
git & 2.50.0 & \texttt{make verify} \\
\bottomrule
\end{tabularx}
\end{table}

Deux outils hors de Python ont servi aux analyses ajoutées après la relecture technique. Le client
\texttt{gh} 2.95.0 a lu la date de création des dépôts GitHub cités pour Q5 (écart 31 du
sondage), une requête à la fois. Ollama 0.40.16 sert localement \texttt{gemma3:12b}, le lecteur de
l'écart 33 du sondage. Chaque fichier de lecture de ce lecteur note le condensé du modèle.

Le présent rapport se construit par \texttt{make report}, qui lance \file{report/build.sh} : les
tableaux depuis les fichiers JSON (\file{report/tools/tables.py}), les figures
(\file{report/tools/figures.py}), puis deux passes de LuaLaTeX avec biber, en français
(\file{rapport.tex}) et en anglais (\file{report-en.tex}).
